\documentclass[twoside,a4paper]{book}

\usepackage[pdftex]{graphicx}
\usepackage{amsmath}
\usepackage{amssymb}
\usepackage{textcomp}
\usepackage[utf8]{inputenc}
\usepackage[polish]{babel}
\usepackage[T1]{fontenc}
\usepackage{array}
% pakiet stosowany do url'i w bibliografii, zamienia odnośniki na ładnie sformatowane
\usepackage{url}
% pakiety służące do numerowania i tworzenia algorytmów
\usepackage{algorithmic}
\usepackage{algorithm}
% redefinicja etykiety nagłówkowej listy algorytmów, domyślna jest po angielsku
\renewcommand{\listalgorithmname}{Spis algorytmów}

% pakiet do wyliczania skali, przydatny przy dużych obrazkach
\usepackage{pgf}
% pakiet służący do automatycznego sortowania odnośników do bibliografii
\usepackage[sort]{natbib}
% tworzenie listingów
\usepackage{listings}
% tworzenie figur wewnątrz figur
\usepackage{subfig}
% do automatycznego skracania nazw rozdziałów i podrozdziałów używanych w nagłówkach strony by mieściły się w jednej linii
\usepackage[fit]{truncate}
% fancyhdr - ładne nagłówki, definicja wyglądu nagłówka, numery stron będą umieszczane w nagłówku po odpowiedniej stronie
\usepackage{fancyhdr}
\pagestyle{fancy}
\renewcommand{\chaptermark}[1]{\markboth{#1}{}}
\renewcommand{\sectionmark}[1]{\markright{\thesection\ #1}}
\fancyhf{}
\fancyhead[LE,RO]{\bfseries\thepage}
% tutaj ograniczamy szerokość pola w nagłówku zawierającego nazwę rozdziału/podrozdziału do 95% szerokości strony
% redefinicja sposobu prezentacji nazw domyślnie wypisywanych wielkimi literami (np. domyślnie w nagłówku Spis treści będzie miał postać SPIS TREŚCI)
% Uwaga! to może popsuć wielkie litery w ogóle! Jak coś nie działa należy usunąć \nouppercase{} z poniższych definicji
\fancyhead[LO]{\nouppercase{\bfseries{\truncate{.95\headwidth}{\rightmark}}}}
\fancyhead[RE]{\nouppercase{\bfseries{\truncate{.95\headwidth}{\leftmark}}}}
\renewcommand{\headrulewidth}{0.5pt}
\renewcommand{\footrulewidth}{0pt}

% definicja typu prostego wymagana przez pierwsze strony rozdziałów itp.
% powyższe reguły niestety tych stron nie dotyczą, gdyż Latex automatycznie przełącza je pomiędzy fancy a plain
% w tym wypadku eliminujemy nagłówki i stopki na stronach początkowych
\fancypagestyle{plain}{%
 \fancyhead{}
 \fancyfoot{}
 \renewcommand{\headrulewidth}{0pt}
 \renewcommand{\footrulewidth}{0pt}
}

\parskip 0.05in


% makro umożliwiające otaczanie symboli okręgami
\usepackage{tikz}
% brak justowania tekstu (bazą okręgu będzie linia tekstu)
\newcommand*\mycirc[1]{%
  \begin{tikzpicture}
    \node[draw,circle,inner sep=1pt] {#1};
  \end{tikzpicture}}

% pionowe justowanie tekstu, środek okręgu pokrywa się ze środkiem tekstu
\newcommand*\mycircalign[1]{%
  \begin{tikzpicture}[baseline=(C.base)]
    \node[draw,circle,inner sep=1pt](C) {#1};
  \end{tikzpicture}}

% zmiana nazwy twierdzeń i lematów
\newtheorem{theorem}{Twierdzenie}[section]
\newtheorem{definition}{Definicja}
\newtheorem{lemma}[theorem]{Lemat}

% tworzenie definicji dowodu
\newenvironment{proof}[1][Dowód]{\begin{trivlist}
\item[\hskip \labelsep {\bfseries #1}]}{\end{trivlist}}
% \newenvironment{definition}[1][Definicja]{\begin{trivlist}
% \item[\hskip \labelsep {\bfseries #1}]}{\end{trivlist}}
% \newenvironment{example}[1][Przykład]{\begin{trivlist}
% \item[\hskip \labelsep {\bfseries #1}]}{\end{trivlist}}
% \newenvironment{remark}[1][Uwaga]{\begin{trivlist}
% \item[\hskip \labelsep {\bfseries #1}]}{\end{trivlist}}

% definicja czarnego prostokąta zwyczajowo dodawanego na koniec dowodu
\newcommand{\qed}{\nobreak \ifvmode \relax \else
      \ifdim\lastskip<1.5em \hskip-\lastskip
      \hskip1.5em plus0em minus0.5em \fi \nobreak
      \vrule height0.75em width0.5em depth0.25em\fi}

% poniższymi instrukcjami można sterować co ma być numerowane a co nie i co ma być wyświetlane w spisie treści
% \setcounter{secnumdepth}{3}
% \setcounter{tocdepth}{5}

% definicja czcionki mniejszej niż tiny (domyślnie takiej małej nie ma)
\usepackage{lmodern}
\makeatletter
  \newcommand\tinyv{\@setfontsize\tinyv{4pt}{6}}
\makeatother

% definicja jeszcze mniejszej czcionki
\usepackage{lmodern}
\makeatletter
  \newcommand\tinyvv{\@setfontsize\tinyvv{3.5pt}{6}}
\makeatother

% pakiet do obsługi wielostronicowych tabel
\usepackage{longtable}
\setlength{\LTcapwidth}{\textwidth}

\usepackage[section] {placeins}

\usepackage{multirow}

\usepackage{slantsc}

% nazwa pliku ze stroną tytułową
% % allows useg @ as a @ not as special character
% required for macro redefinition
\makeatletter

% parameters definition
% they cannot conflict with other
% like bibteh attributes etc.
\def\promotor#1{\def\@promotor{#1}}
\def\promotordpt#1{\def\@promotordept{#1}}
\def\promotoruniv#1{\def\@promotoruniv{#1}}
\def\miasto#1{\def\@miasto{#1}}

\def\maketitle{
  %removal of header
  \thispagestyle{empty}%
  %changing margins to match department requirements
%   \changepage{+60pt}{+70pt}{}{-30pt}{}{-30pt}{-12pt}{-18pt}{}
  \begin{center}
    \begin{tabular}{@{}lcr@{}}
      \multirow{4}{*}{\includegraphics[height=1.9cm]{img/logo_pg.png}} &
      \LARGE{\textbf{POLITECHNIKA GDAŃSKA}} &
      \multirow{4}{*}{\includegraphics[height=1.9cm]{img/logo_eti.png}}\\
      & & \\
      &\LARGE{\textbf{Wydział Elektroniki,}}&\\
      &\LARGE{\textbf{Telekomunikacji i Informatyki}}&
    \end{tabular}
  \end{center}
  \rule{\linewidth}{0.1mm}
  \vspace{2.82cm}
  \begin{center}
    \huge{\textbf{\@author}}
  \end{center}
  \vspace{0.5cm}
  \begin{center}
    \Huge{\textbf{\@title}}
  \end{center}
  \vspace{0.5cm}
  \begin{center}
    \huge{Rozprawa doktorska}
  \end{center}
  \vspace{1.41cm}
% \renewcommand{\arraystretch}{0.85}
  \begin{tabular}{@{}p{0.46\textwidth}@{}@{}p{0.54\textwidth}@{}}
\noalign{\smallskip}
    & \hspace{-20pt}\LARGE{Promotor:}\\
\noalign{\smallskip}
\noalign{\smallskip}
    & \Large{\@promotor}\\
\noalign{\smallskip}
\noalign{\smallskip}
    & \Large{\@promotordept}\\
\noalign{\smallskip}
\noalign{\smallskip}
    & \Large{\@promotoruniv}
  \end{tabular}
  \vfill
  \begin{center}
    \Large{\@miasto, \@date}
  \end{center}
  %
%   \cleardoublepage
  % restore page defaults
%   \changepage{-60pt}{-70pt}{}{+30pt}{}{+30pt}{+12pt}{+18pt}{}
}

%restore @ sign
\makeatother

\cleardoublepage
% allows useg @ as a @ not as special character
% required for macro redefinition
\makeatletter

% parameters definition
% they cannot conflict with other
% like bibteh attributes etc.
\def\promotor#1{\def\@promotor{#1}}
\def\miasto#1{\def\@miasto{#1}}
\def\studies#1{\def\@studies{#1}}
\def\descr#1{\def\@descr{#1}}
\def\indeks#1{\def\@indeks{#1}}
\def\dept#1{\def\@dept{#1}}

\def\maketitle{
  %removal of header
  \thispagestyle{empty}%

  \begin{center}
    \begin{tabular}{lcl}
      \multirow{4}{*}{\includegraphics[height=2.5cm]{img/logo_pg.png}} &
      \textsc{\textbf{Politechnika Gdańska}} &
      \multirow{4}{*}{\includegraphics[height=2.5cm]{img/logo_eti.png}}\\
      & & \\
      &\textsc{\textbf{Wydział Elektroniki,}}&\\
      &\textsc{\textbf{Telekomunikacji i Informatyki}}&
    \end{tabular}
  \end{center}
  \vspace{1cm}
  \begin{tabular*}{\textwidth}{p{0.5\textwidth}p{0.5\textwidth}}
    \textbf{Katedra:} & \@dept\\
    &\\
    \textbf{Imię i nazwisko dyplomanta:} & \@author\\
    &\\
    \textbf{Nr albumu:} & \@indeks\\
    &\\
    \textbf{Forma i poziom studiów:} & \@studies\\
    &\\
    \textbf{Kierunek studiów:} & Informatyka\\
  \end{tabular*}
  \begin{center}
    \vspace{1cm}
    \Large{\textbf{Praca dyplomowa magisterska}}
  \end{center}
  \vspace{1cm}
  \begin{tabular*}{\textwidth}{p{\textwidth}}
    \textbf{Temat pracy:} \\ \@title\\
    \\
    \textbf{Kierujący pracą:} \\ \@promotor\\
    \\
    \textbf{Zakres pracy:} \\ \@descr\\
  \end{tabular*}
  \vspace*{\stretch{6}}
  \begin{center}
    \@miasto, \@date
  \end{center}

}

%restore @ sign
\makeatother

\cleardoublepage

% parametry strony tytułowej, zdefiniowane są w plikach z poszczególnymi stronami
% tytuł pracy
\title{Szablon pracy magisterskiej i~doktorskiej, pamiętaj by odkomentować właściwy plik ze stroną tytułową!}
% autor
\author{Tomasz Boiński}
% rok wydania
\date{2012}
% miasto, gdzie napisano pracę
\miasto{Gdańsk}
% promotor
\promotor{prof.\ dr hab.\ inż.\ Jan Iksiński}
% wydział promotora, tylko dla phd_titlepage
% \promotordpt{Wydział Elektroniki, Telekomunikacji i~Informatyki}
% uczelnia promotora, tylko dla phd_titlepage
% \promotoruniv{Politechnika Gdańska}

% rodzaj studiów, tylko dla mgr_titlepage
\studies{Stacjonarne jednolite studia magisterskie}
% opis pracy, tylko dla mgr_titlepage
\descr{Opis pracy (jedno zdanie)}
% nr indeksu, tylko dla mgr_titlepage
\indeks{999999}
% katedra, tylko dla mgr_titlepage
\dept{Architektury Systemów Komputerowych}

% korekta marginesów - domyślnie latex ma jakieś kosmiczne
\usepackage{anysize}
\marginsize{3.5cm}{2.5cm}{2.5cm}{2.5cm}
% po zmianie marginesów konieczne jest wymuszenie przeliczenia nagłówków
\fancyhfoffset[E,O]{0pt}

\begin{document}
% sekcja wstępna książki, numerowana rzymskimi
\frontmatter
% generacja strony tytułowej załączonej wcześniej
\maketitle
% spis treści
\tableofcontents

% dodatkowa strona z podaniem źródła finansowania itp.
\newpage
% wstawienie pustego symbolu
\null
% i wypełnienie nim całej dostępnej strony
\vfill
% a na jej dole dodajemy właściwy tekst
% w tym przypadku z wyrównaniem do prawej strony
\begin{flushright}
  Wszystko co zostało dodane we frontmatter zostanie wydzielone do odrębnie numerowanej sekcji. Numeracja zazwyczaj jest małymi, rzymskimi liczbami. Tutaj przykład jak tekst umieścić na dole strony z wyrównaniem do prawej. Inne układy można robić analogicznie.
\end{flushright}

% dodatkowy, nienumerowany rozdział na np. podziękowania
\chapter*{Podziękowania}
Tu możemy dać podziękowania. Jest to zrealizowane w~postaci zwykłego, nienumerowanego rozdziału. W razie konieczności należy ten rozdział po prostu usnąć.















% właściwa część książki, numerowana arabskimi od 1
\mainmatter

\chapter{Przegląd literatury}
\section{Wielomiany}
\subsection{Definicja wielomianu}




\begin{definition}
	Wielomianem zmiennej rzeczywistej x nazywamy wyrażenie:
	
	\begin{equation}
	\begin{split}
	&W(x) = a_0x^m + a_1x^{m-1} a_2x+ ... + a_{m-1}x + a_m, \\
	&gdzie\ liczby\ a_0, a_1, a_2, ..., 1_{m-1}, a_m\ są liczbami\ rzeczywistymi
	\end{split}
	\label{eq:lin:defPolynomialeq}
	\end{equation}
	\label{eq:lin:defPolynomial}
\end{definition}













\subsection{Twierdzenia i dowody}

Linie 75 -- 93 nagłówka dokumentu definiują nowe nazwy sekcji twierdzeń i dowodów, oraz znacznik końca dowodu (taki czarny kwadracik). Dzięki nim można uzyskać ładnie wyglądające twierdzenia jak poniżej (Twierdzenie~\ref{eq:lin:theorem}). Zauważmy, że równanie dowodu nie jest równaniem numerowanym. Wszędzie tam, gdzie nie chcemy by rozdział czy dowolna inna sekcja była numerowana należy w jej nazwie użyć gwiazdki, np. \textbf{$\backslash$begin\{equation*\}}.

\begin{theorem}
	Podobieństwo pomiędzy pojęciami $A$ i~$B$ opisane jest stosunkiem ilości informacji niezbędnej do opisania ich wspólności znaczeniowej oraz ilością informacji niezbędnej do ich opisania (Równanie~\ref{eq:lin:theoremeq}).
	
	\begin{equation}
	sim_{lin}(A,B)=\frac{\log P(common(A,B))}{\log P(description(A,B))}
	\label{eq:lin:theoremeq}
	\end{equation}
	\label{eq:lin:theorem}
\end{theorem}

\begin{proof}
	\begin{equation*}
	\begin{split}
	f(x,y)&=f(x+0,y+(y-x))\\
	&=\frac{x}{y}*f(x,x) +~\frac{y-z}{x}*f(0,y-z)\\
	&=\frac{x}{y}*1 +~\frac{y-z}{x}*0\\
	&=\frac{x}{y} \qed
	\end{split}
	\end{equation*}
\end{proof}

Inna ciekawa konstrukcja wykorzystująca tryb matematyczny do zapisania pewnego stwierdzenia:

Niech $A \subseteq T$, $C = N_y(A) \neq W$, a~$\alpha_y = \min_{a \in A, b\notin C} \{y(a) +~y(b) - q(a, b)\}$ oraz

\[ y'(v) = \left\{ \begin{array}{ll}
y(v) - \alpha_y & \mbox{jeżeli $v \in A$} \\
y(v) +~\alpha_y & \mbox{jeżeli $v \in C$} \\
y(v)            & \mbox{w innych przypadkach}
\end{array}
\right. \]

Zapis ten, acz skomplikowany, pozwala na reprezentację złożonych reguł matematycznych w postaci ładnie ułożonych i wyrównanych wierszy. Reguły \textbf{left} oraz \textbf{right} pozwalają na utworzenie nawiasów klamrowych, których rozmiar będzie automatycznie dostosowywany do rozmiaru elementu, jakie mają zawierać.


















\chapter{Wstęp}
\section{Idea dokumentu}

Dokument ten ma w~założeniu pomóc w pisaniu pracy magisterskiej/rozprawy doktorskiej w narzędziu jakim jest Latex. Zdecydowanie polecam jego używanie, gdyż znacznie upraszcza formatowanie i~skład tekstu, wykonując bardzo dużo pracy za autora. W dalszej części tego dokumentu na przykładach przedstawię kilka technik jak zrealizować rzeczy, na jakie natknąłem się w~trakcie pisania mojej pracy magisterskiej oraz rozprawy doktorskiej.

Dokument ten nie ma na celu dostarczenie kompletnej dokumentacji Latex'a, czy też zaznajomić z nim osoby nie mającego żadnej wiedzy na jego temat. Nie jest to samouczek (takich jest cała masa w sieci, np.~\cite{wikibooks}, polecam też zapoznanie się z dokumentacją pakietów, jakie zostały użyte w tym dokumencie), ale raczej zbiór ciekawych konstrukcji, jakie znalazłem w~trakcie pisania swojej pracy magisterskiej i rozprawy doktorskiej. Jednakże tam gdzie uznam to za stosowne będę podawał rzeczy również oczywiste, by osoby nie mające wcześniej kontaktu z Latex'em mogły, na bazie tego opracowania, rozpocząć pracę nad swoim dokumentem. Domyślnie przygotowany jest dla dwustronnego wydruku w stylu książki. Wydruk jednostronny wymaga zmiany definicji klasy z \textbf{twoside} na \textbf{oneside} i poprawienia definicji nagłówków i stopek (linie 35 -- 56 preambuły dokumentu). Polecam również zaglądać do kodu dokumentu. Sam dokument bardziej pokazywać będzie efekt końcowy, jaki można uzyskać. Samą treść najłatwiej podejrzeć już w pliku źródłowym.

Do pisania samego tekstu użyć można dowolnego edytora tekstu, pozwalającego na tworzenie zwykłych plików tekstowych. Polecam jednak jakieś zintegrowane środowisko umożliwiające szybką kompilację i~podgląd wynikowego dokumentu. Osobiście używam narzędzia o nazwie \textbf{Kile} dostępnego praktycznie w każdej dystrybucji systemu Linux.

Zalecam przyjrzenie się również nagłówkowi dokumentu -- zawiera on importy różnych przydatnych pakietów dodatkowych wraz z~opisem. Większość z nich będzie opisana lub wykorzystana w tym dokumencie.

Gorąco zachęcam do korzystania z Latex'a. Przy odrobinie wprawy jest to narzędzie wysoce wygodne i~wydajne, umożliwiające zdecydowanie prostsze tworzenie złożonych dokumentów niż dowolny WYSYWIG\@.

\section{Strony tytułowe}

Do dokumentu dołączyłem strony tytułowe podobne do oficjalnych formatek pracy magisterskiej (mgr\_titlepage.tex) oraz rozprawy doktorskiej(phd\_titlepage.tex). Pierwsza z nich, będąca bazą dla drugiej, opracowana została przez mgra inż.\ Michała Wójcika\footnote{\url{http://mwojcik.eti.pg.gda.pl}}.

\section{Poprawność dokumentu}

Przed drukiem warto zawsze sprawdzić swój dokument pod względem formalnej poprawności. Służy do tego narzędzie \textbf{lacheck}. Na stronach Katedry Architektury Systemów Komputerowych znajduje się interfejs WWW do tego narzędzia napisany przez mgra inż.\ Rafała Knopę. Narzędzie to jest dostępne pod adresem \url{http://kask.eti.pg.gda.pl/lacheck/}.

\section{Uwagi i poprawki}

Zachęcam wszystkich do modyfikacji, poprawiania i rozbudowy niniejszego dokumentu. Z Waszą pomocą z biegiem czasu dokument ten ma szansę stać się kompletnym przewodnikiem do Latexa pomocnym w trakcie pisania pracy magisterskiej czy doktorskiej. Wszelkie zmiany przesłane do mnie opublikuję na stronie katedralnej, oczywiście z zaznaczeniem współautorów.

\chapter{Podstawowe elementy}

To jest rozdział.

\section{Sekcje dokumentu}

W Latexu w klasie dokumentów \textbf{book} wyróżniamy rozdziały (\textbf{chapter}), podrozdziały \textbf{section}, podpodrozdziały \textbf{subsection}, podpodpodrozdziały \textbf{subsubsection} i paragrafy (\textbf{paragraph}). Podpodpodrozdziały i paragrafy domyślnie nie są numerowane ani nie występują w spisie treści. Zachowanie to można zmienić poprzez funkcję \textbf{setcounter} umieszczaną w preambule. Wykomentowany przykład można znaleźć w kodzie tego dokumentu.

Obecnie znajdujemy się na poziomie podrozdziału. Pozostałe przykłady poniżej.

\subsection{Podpodrozdział}

To jest podpodrozdział.

\subsubsection{Podpodpodrozdział}

To jest podpodpodrozdział. On nie jest domyślnie numerowany.

\paragraph{Paragraf}

A to jest paragraf. On również nie jest domyślnie numerowany.

\section{Podstawowe elementy typograficzne}

\subsection{Twarda spacja}

Twarda spacja jest bardzo istotnym elementem, gdyż zabrania Latex'owi łamanie linii w miejscu jej wystąpienia, a tym samym pozwoli na niejako ,,sklejenie'' wyrazów ze sobą. Dzięki temu możemy uniknąć tzw.\ sierot (pojedynczych znaków na końcu wiersza). W Latex twardą spację umieszcza się wstawiając znak tyldy~(\textasciitilde). Zapisujemy to więc np.\ tak: ,,dokument, w{\textasciitilde}którym''.

\subsection{Formatowanie tekstu}

Aby zapewnić poprawny wygląd tekstu należy pamiętać o kilku rzeczach:

\begin{itemize}
 \item Linia poprzedzona procentem to komentarz.
 \item Poprzedzaniu spacji występującej po kropce kończącej skrót znakiem ucieczki, odstęp będzie wtedy taki, jak odstęp między wyrazami a~nie między zdaniami. Przykładowo zapis ,,np. tekst'' vs. ,,np.\ tekst''. Ten drugi jest poprawny, a zapisany został tak: ,,np.$\backslash$~tekst''.
 \item Skróty pisane wielkimi literami kończące zdanie powinny posiadać {$\backslash$}@ przed kropką kończącą zdanie, np.\ OCS{$\backslash$}@\@. Spowoduje to potraktowanie spacji jako spacji międzyzdaniowej z nie międzywyrazowej.
 \item Cudzysłowie zawsze tworzymy używając podwójnego przecinka jako symbolu otwierającego cudzysłów, oraz podwójnego apostrofu zamykającego cudzysłów.
 \item Kursywę uzyskujemy za pomocą słowa kluczowego {$\backslash$}textit, co w efekcie daje \textit{tekst kursywą}. Pogrubiony \textbf{używamy słowa kluczowego textbf}. Każdorazowo tekst mający być napisany danym krojem otaczamy nawiasami klamrowymi.
 \item Myślnik (--) tworzymy poprzez umieszczenie bezpośrednio po sobie dwu kresek (minusy). Różnica między nimi jest zasadnicza. Pojedynczy myślnik generuje krótką kreskę (-), podwójny długą (--), potrójny najdłuższą (---).
 \item Odwołania do różnych elementów dokumentu robimy poprzez słowo kluczowe \textbf{ref()}. Jako jego parametr wstawiamy nazwę zdefiniowaną za pomocą słowa kluczowego \textbf{label()}. Należy pamiętać, że odwołanie zwraca jedynie numer elementu, słowo opisowe, jak np.\ rozdział czy rysunek należy dodać samodzielnie. Polecam tutaj przyjąć jakąś konwencję i się jej trzymać w całym dokumencie. Tak samo należy postępować w przypadku etykiet.
 \item Latex doskonale radzi sobie z dzieleniem wyrazów na końcach linii, jednak czasami zachodzi konieczność wymuszenia podziału w określonym miejscu. W tym celu należy zastosować konstrukcję $\backslash$-. Latex takiego ukośnika nie wydrukuje dopóty, dopóki rzeczywiście w tym miejscu nie zostanie wykonane przeniesienie części wyrazu. Możliwe jest dodanie wielu podziałów w jednym wyrazie. Użyte wtedy zostanie to, które spowoduje wygenerowania ,,najładniejszego'' tekstu.
\end{itemize}

\section{Podział linii i paragrafy}
\label{podzial}

Nowy paragraf rozpoczyna się poprzez wstawienie jednej wolnej linii. Latex automatycznie wygeneruje wcięcie. Należy pamiętać, że pierwszy paragraf, zgodnie ze standardami drukarskimi, nie ma wcięcia! Możemy tym sterować za pomocą poleceń \textbf{noindent} (brak wcięcia) oraz \textbf{indent} (dodatkowe wcięcie).

Jeżeli chcemy po prostu zrobić nową linię, bez tworzenia nowego paragrafu używamy konstrukcji $\backslash\backslash$. Efekt będzie taki, że paragraf\\
będzie kontynuowany w nowej linii. Nie spowoduje to jednak rozciągnięcia poprzedniej linii. Zostanie ona przerwana tam gdzie tego sobie zażyczymy i kontynuowana w nowej linijce.

A co w przypadku, gdy chcemy z jakiegoś powodu przerwać linię, ale wymusić justowanie tekstu? Weźmy dla przykładu fragment:

Trzecim istotnym aspektem jest stosowana w~trakcie wytwarzania ontologii metodologia pracy~\cite{boinski2012kaskbook,boinski2011security}. Zastosowanie jednej z~uznawanych metodologii, takich jak Methontology, NeOn czy metodologia opracowana przez Noy~i~McGuiness, znacząco wpływa na jakoś uzyskanego produktu. Wspomniane metodologie w~dużej mierze uwzględniają potrzebę przyszłej integracji wiedzy, a~w połączeniu z~narzędziami typu Protégé czy OCS~\cite{boinski2007kaskbook,boinski2009ocs,boinski2010zespolowa}, pozwalają na tworzenie spójnych i~formalnie oraz logicznie poprawnych ontologii.

Tekst zostaje bardzo brzydko złamany w środku odnośników do cytowań. Użycie podwójnego po słowie ,,metodologia'' w pierwszym zdaniu ukośnika da nam natomiast taki efekt:

Trzecim istotnym aspektem jest stosowana w~trakcie wytwarzania ontologii metodologia \\ pracy~\cite{boinski2012kaskbook,boinski2011security}. Zastosowanie jednej z~uznawanych metodologii, takich jak Methontology, NeOn czy metodologia opracowana przez Noy~i~McGuiness, znacząco wpływa na jakoś uzyskanego produktu. Wspomniane metodologie w~dużej mierze uwzględniają potrzebę przyszłej integracji wiedzy, a~w połączeniu z~narzędziami typu Protégé czy OCS~\cite{boinski2007kaskbook,boinski2009ocs,boinski2010zespolowa}, pozwalają na tworzenie spójnych i~formalnie oraz logicznie poprawnych ontologii.

Też nie ładnie, gdyż linijka jest niewyjustowana. Z pomocą przychodzi nam tutaj komenda \textbf{linebreak[]}, gdzie w nawiasie kwadratowym podajemy liczbę od 1 do 4 określająca jak bardzo zależy nam na tym, by linia została złamana w tym miejscu (4 to najwyższa wartość). Efekt jest następujący:

Trzecim istotnym aspektem jest stosowana w~trakcie wytwarzania ontologii metodologia \linebreak[4] pracy~\cite{boinski2012kaskbook,boinski2011security}. Zastosowanie jednej z~uznawanych metodologii, takich jak Methontology, NeOn czy metodologia opracowana przez Noy~i~McGuiness, znacząco wpływa na jakoś uzyskanego produktu. Wspomniane metodologie w~dużej mierze uwzględniają potrzebę przyszłej integracji wiedzy, a~w połączeniu z~narzędziami typu Protégé czy OCS~\cite{boinski2007kaskbook,boinski2009ocs,boinski2010zespolowa}, pozwalają na tworzenie spójnych i~formalnie oraz logicznie poprawnych ontologii.

Jeżeli z jakiegoś powodu potrzebujemy nową linię to używamy komendy \textbf{newpage}. \newpage Tekst występujący po niej znajdzie się na nowej stronie. Rozdziały itp.\ automatycznie generują nową stronę, przy czym w układzie dwustronnym nowy rozdział zawsze zacznie się od nieparzystej strony.

\section{Środowisko matematyczne}

Środowisko matematyczne otwieramy i zamykamy znakiem \$. Niektóre funkcje można używać tylko wewnątrz takiego środowiska. Przykładem niech będzie funkcja \textbf{mathcal} zamieniająca duże litery w symbole o charakterystycznym kroju, stosowanym do opisywania stałych, np.\ $\mathcal{O}$ czy $\mathcal{R(D,P,T,S,U,I)}$. Pamiętać należy, że zamienione zostaną wszystkie litery w wyrażeniu występującym wewnątrz nawiasów klamrowych.

Niektóre konstrukcje, np.\ równania, automatycznie włączają tryb matematyczny. Równania dobrze jest opisać, przykład przedstawia Równanie~\ref{eq:przyklad}.

\begin{equation}
  \mathcal{O(K,B,C,R)}
  \label{eq:przyklad}
\end{equation}

gdzie:\\
$\mathcal{K}$ - zbiór klas wchodzących w~skład ontologii,\\
$\mathcal{B}$ - zbiór bytów wchodzących w~skład ontologii,\\
$\mathcal{C}$ - zbiór komentarzy przypisanych do klas $\mathcal{K}$ i~bytów $\mathcal{B}$ wchodzących w~skład ontologii,\\
$\mathcal{R}$ - zbiór relacji wiążących elementy ontologii.


W równaniach możemy stosować różne dodatkowe symbole oraz np.\ wyrównywać je do określonego miejsca. Służy do tego blok typu \textbf{split}, a sam punkt wyrównania określony jest ampersandem (\&). Przykład zastosowania prezentuje równanie~\ref{eq:split_ex}.

\begin{equation}
  \label{eq:split_ex}
  \begin{split}
    \forall {x_1 \leq y_1, x_2 \leq y_2}&: f(x_1+x_2,y_1+y_2)\\ 
				  &= \frac{y_1}{y_1+y_2}f(x_1,y_1)+\frac{y_2}{y_1+y_2}f(x_2,y_2)
  \end{split}
\end{equation}

\subsection{Twierdzenia i dowody}

Linie 75 -- 93 nagłówka dokumentu definiują nowe nazwy sekcji twierdzeń i dowodów, oraz znacznik końca dowodu (taki czarny kwadracik). Dzięki nim można uzyskać ładnie wyglądające twierdzenia jak poniżej (Twierdzenie~\ref{eq:lin:theorem}). Zauważmy, że równanie dowodu nie jest równaniem numerowanym. Wszędzie tam, gdzie nie chcemy by rozdział czy dowolna inna sekcja była numerowana należy w jej nazwie użyć gwiazdki, np. \textbf{$\backslash$begin\{equation*\}}.

\begin{theorem}
 Podobieństwo pomiędzy pojęciami $A$ i~$B$ opisane jest stosunkiem ilości informacji niezbędnej do opisania ich wspólności znaczeniowej oraz ilością informacji niezbędnej do ich opisania (Równanie~\ref{eq:lin:theoremeq}).

 \begin{equation}
   sim_{lin}(A,B)=\frac{\log P(common(A,B))}{\log P(description(A,B))}
   \label{eq:lin:theoremeq}
 \end{equation}
 \label{eq:lin:theorem}
\end{theorem}

\begin{proof}
  \begin{equation*}
   \begin{split}
     f(x,y)&=f(x+0,y+(y-x))\\
	   &=\frac{x}{y}*f(x,x) +~\frac{y-z}{x}*f(0,y-z)\\
	   &=\frac{x}{y}*1 +~\frac{y-z}{x}*0\\
	   &=\frac{x}{y} \qed
   \end{split}
 \end{equation*}
\end{proof}

Inna ciekawa konstrukcja wykorzystująca tryb matematyczny do zapisania pewnego stwierdzenia:

Niech $A \subseteq T$, $C = N_y(A) \neq W$, a~$\alpha_y = \min_{a \in A, b\notin C} \{y(a) +~y(b) - q(a, b)\}$ oraz

\[ y'(v) = \left\{ \begin{array}{ll}
                y(v) - \alpha_y & \mbox{jeżeli $v \in A$} \\
                y(v) +~\alpha_y & \mbox{jeżeli $v \in C$} \\
                y(v)            & \mbox{w innych przypadkach}
               \end{array}
       \right. \]

Zapis ten, acz skomplikowany, pozwala na reprezentację złożonych reguł matematycznych w postaci ładnie ułożonych i wyrównanych wierszy. Reguły \textbf{left} oraz \textbf{right} pozwalają na utworzenie nawiasów klamrowych, których rozmiar będzie automatycznie dostosowywany do rozmiaru elementu, jakie mają zawierać.

chapter{Rysunki, tabele i inne konstrukcje}
\enlargethispage{20pt}
section{Rysunki}
\label{rysunki}
Obrazki wstawiamy do dokumentu za pośrednictwem sekcji \textbf{figure}. Przykładowy wyśrodkowany rysunek z etykietą i opisem to Rys.~\ref{fig:obrazek}.

\begin{figure}[!h]
\centering
\includegraphics[width=0.7\textwidth]{img/obrazek.png}
\caption{Opis do obrazka.}
\label{fig:obrazek}
\end{figure}

\section{Algorytmy}
\label{algorytmy}
Pakiety \textbf{algorithmic} oraz \textbf{algorithm} dostarczają środowiska do definiowania algorytmów. Można za ich pomocą definiować zarówno pseudokod jak i algorytmy ogólne. Pseudokod przedstawia Algorytm~\ref{tboi:alg}. Przykładową ogólną reprezentację algorytmu Levenshteina obrazuje Algorytm~\ref{tboi:alg_leven}.

\begin{algorithm}
 \floatname{algorithm}{Algorytm}
 \caption{Pseudokod prezentujący jakiś bezsensowny algorytm}
 \label{tboi:alg}
 \begin{algorithmic}[1]
\STATE \textbf{program} FancyProgram \{

\STATE \textbf{function} superFunkcja(List listA) \{

\FORALL{element :~listA.getElements()}
\STATE int result = storeElement(element);
\IF{result == SUCCESS}
\STATE double number = result*element;
\ELSIF{result == FAILURE}
\STATE double number = result/element;
\ENDIF
\ENDFOR
\STATE \}

\STATE input List elements;
\STATE output number;

\STATE superFunkcja(elements);
\RETURN number;
\STATE \}
 \end{algorithmic}
\end{algorithm}

\begin{algorithm}
 \floatname{algorithm}{Algorytm}
 \caption{Algorytm Levenshteina}
 \label{tboi:alg_leven}
%  \begin{algorithmic}[1]
\begin{enumerate}
 \item Niech $n = \textit{długość}(s)$ a~$m = \textit{długość}(t)$.
 \item Jeżeli $n == 0$ zwróć m~i zakończ działanie.
 \item Jeżeli $m == 0$ zwróć n~i zakończ działanie.
 \item Utwórz macierz $d$ o~$m+1$ wierszach indeksowanych od $0$ do $m$ oraz $n+1$ kolumnach indeksowanych od $0$ do $n$, pierwszy wiersz zainicjuj wartościami od $0$ do $n$ a~pierwszą kolumnę od $0$ do $m$.
 \item Porównaj każdy znak pierwszego ciągu, indeksowany za pomocą $i$, z~każdym znakiem drugiego ciągu, indeksowanym za pomocą $j$. Dla każdego porównania:
\begin{enumerate}
 \item jeżeli $s[i]$ jest identyczne z~$t[j]$, $koszt = 0$, w~przeciwnym wypadku $koszt = 1$,
 \item ustaw wartość komórki $d[i,j]$ macierzy na wartość minimalną spośród:
\begin{itemize}
 \item $d[i-1,j] +~1$,
 \item $d[i,j-1] +~1$,
 \item $d[i-1,j-1] +~koszt$.
\end{itemize}
\end{enumerate}
 \item Odczytaj odległość $dist_{lev}(s,t)$ z~komórki $d[n,m]$.
\end{enumerate}
%\end{algorithmic}
\end{algorithm}

Podobnie jak w przypadku rysunków, algorytmy są automatycznie rozlokowywane w tekście przez system Latex.


\section{Wpisy bibliograficzne}

Wpisy bibliograficzne przechowujemy w odrębnym pliku z rozszerzeniem\ .bib. Przykładowy plik został dołączony do tego dokumentu. Do pliku takiego należy dodawać odpowiedni sformatowane wpisy. Latex automatycznie posortuje je po nazwiskach autorów oraz do finalnego dokumentu dołączy tylko te wpisy,które posiadają odwołania w tekście! Można więc stworzyć kompletną bazę publikacji, a Latex użyje tylko to co potrzeba. Dodatkowo, dzięki użyciu pakietu \textbf{natbib} z parametrem \textbf{sort} (patrz preambuła dokumentu), numerki w odwołaniach również zostaną posortowane, niezależnie od kolejności podania odwołań. Przykład podano w rozdziale~\ref{podzial}.


\backmatter

% lista rysunków
\listoffigures
% lista tabel
\listoftables
% lista algorytmów
\listofalgorithms

% rodzaj bibliografii
\bibliographystyle{plain}
% plik z wpisami bibliograficznymi
\bibliography{bibliografia}
\end{document}